\BourbakiExerciseGroup

\begin{enumerate}
\def\labelenumi{\arabic{enumi})}
\tightlist
\item
  （指数型序列的代数）给定两个整数 \(p, q \in \mathbf{N}\) ，我们用 \(((p, q))\) 表示二项式系数 \(\frac{(p + q)!}{p!q!}\) 。设 \(k\) 为一个交换环， \(B\) 为一个结合的交换含幺 \(k\) -代数。所谓 \(B\) 中元素的指数型序列，我们指的是这样的序列 \(a = (a_p)_{p \in \mathbf{N}}\) ：\(a_0 = 1_B\) 且 \(a_p a_q = ((p, q)) a_{p + q}\) 对任意 \(p, q \in \mathbf{N}\) 成立。 \(B\) 的所有指数型序列的集合记作 \(\mathcal{E}(B)\) 。 \(a)\) 验证在 \(\mathcal{E}(B)\) 上存在唯一的 \(k\) -代数结构，使得对任意序列
\end{enumerate}

\[
\begin{array}{c} (a + b) _ {p} = \sum_ {r + s = p} a _ {r} b _ {s}, \\ (a b) _ {p} = p! a _ {p} b _ {p}, \\ (\lambda a) _ {p} = \lambda^ {p} a _ {p}, \end{array}
\]

对任意序列 \(a, b \in \mathcal{E}(\mathbf{B})\) 以及 \(p \in \mathbf{N}, \lambda \in k\) 成立。映射 \(a \mapsto a_1\) 是代数 \(\mathcal{E}(\mathbf{B})\) 到 \(\mathbf{B}\) 中的同态。

\begin{enumerate}
\def\labelenumi{\alph{enumi})}
\setcounter{enumi}{1}
\item
  假设 \(k\) 包含一个同构于 \(\mathbf{Q}\) 的子环。证明对每个 \(x \in B\)，序列 \(f(x)\) （由 \(f(x)_p = \frac{1}{p!} x^p\) 定义）是指数型的。证明 \(f\) 是 \(B\) 到 \(\mathcal{E}(B)\) 上的一个同构。
\item
  设 \(\mathbf{B}'\) 是一个结合交换含幺 \(k\) -代数，并设 \(h\) 是 \(\mathbf{B}\) 到 \(\mathbf{B}'\) 中的一个同态。对每个 \(a \in \mathcal{E}(\mathbf{B})\)，序列 \(a'\) （由 \(a_p' = h(a_p)\) 对任意 \(p \in \mathbf{N}\) 定义）属于 \(\mathcal{E}(\mathbf{B}')\) 。证明映射 \(\mathcal{E}(h): a \mapsto a'\) 是 \(\mathcal{E}(\mathbf{B})\) 到 \(\mathcal{E}(\mathbf{B}')\) 中的同态。验证 \(\mathcal{E}(g \circ h) = \mathcal{E}(g) \circ \mathcal{E}(h)\) ，其中 \(g\) 和 \(h\) 是可复合的 \(k\) -代数同态。
\item
  设 E 是一个 \(k\) -模，并且 \(\mathbf{TS}(\mathbf{E})\) 是 E 上对称张量的代数。证明对每个 \(x \in \mathbf{E}\) ， \((\gamma_p(x))\) 是代数 \(\mathbf{TS}(\mathbf{E})\) 中的指数型序列（IV，第 45 页）。
\end{enumerate}

\begin{enumerate}
\def\labelenumi{\arabic{enumi})}
\setcounter{enumi}{1}
\tightlist
\item
  （函子 \(\Gamma\) ）记 \(k\) 为一个交换环，E 为一个 \(k\) -模。所谓 E 的伽马代数，我们指的是由生成系 \(\mathbf{N} \times \mathbf{E}\) 及关系元（参见 III，第 450 页）
\end{enumerate}

\[
\begin{array}{c} (0, x) - 1, \\ (p, \lambda x) - \lambda^ {p} (p, x), \\ (p, x + y) - \sum_ {r + s = p} (r, x) (s, y), \\ (p, x) (q, x) - ((p, q)) (p + q, x), \end{array}
\]

，其中 \(p, q\) 遍历 \(\mathbf{N}, x, y\) 遍历 \(\mathbf{E}\)，而 \(\lambda\) 遍历 \(k\)，所定义的结合、交换、含幺代数。该代数记作 \(\Gamma(\mathbf{E})\) 。对每个 \(p \in \mathbf{N}\)，我们用 \(\gamma_p\) 表示 \(\mathbf{E}\) 到 \(\Gamma(\mathbf{E})\) 中的映射，它由单射 \(x \mapsto (p, x)\) 与 \(\mathbf{N} \times \mathbf{E}\) 上的自由交换代数到 \(\Gamma(\mathbf{E})\) 的典范同态复合而成。

\begin{enumerate}
\def\labelenumi{\alph{enumi})}
\item
  验证对每个 \(x \in \mathbb{E}\) ，序列 \(\gamma_p(x)\) 是 \(\Gamma(\mathbb{E})\) 中元素的指数型序列（练习 1），且映射 \(\gamma : x \mapsto (\gamma_p(x))_{p \in \mathbb{N}}\) 是 \(k\) -模 \(\mathbb{E}\) 到 \(k\) -代数 \(\mathcal{E}(\Gamma(\mathbb{E}))\) 中的同态。
\item
  （\(\Gamma(\mathrm{E})\) 的泛性质）。设 B 为一个结合交换含幺 k-代数，且 \(\varphi\) 为 k-模 E 到 \(\mathcal{E}(\mathrm{B})\) 中的同态。证明存在唯一的代数同态 \(h: \Gamma(\mathrm{E}) \to \mathrm{B}\) 使得 \(\varphi = \mathcal{E}(h) \circ \gamma\) 。
\item
  证明存在唯一的含单位元的代数同态 \(\varepsilon : \Gamma(E) \to k\) 使得 \(\varepsilon(\gamma_p(x)) = 0\) 对所有 \(p > 0\) 以及所有 \(x \in E\) 成立。
\end{enumerate}

\(d)\) 对每个 \(p \in \mathbb{N}\) ，令 \(\Gamma_p(E)\) 为 \(\Gamma(E)\) 中由下列元素生成的子模

\[
\gamma_ {\nu_ {1}} (x _ {1}) \dots \gamma_ {\nu_ {s}} (x _ {s})
\]

，其中 \(s \in \mathbf{N}\) ， \(\nu \in \mathbf{N}^s\) 且 \(|\nu| = \sum_{i} \nu_i = p\) 。验证子模 \(\Gamma_p(\mathrm{E})\) 构成代数

\(\Gamma(E)\) 的一个分次。证明 \(\gamma_{1}\) 是 \(E\) 到 \(\Gamma_{1}(E)\) 上的同构（为验证 \(\gamma_{1}\) 是单射，考虑代数 \(B = k \times E\) ，其乘法由公式 \((\lambda, x)(\mu, y) = (\lambda \mu, \lambda y + \mu x)\) 定义，并证明存在 \(\Gamma(E)\) 到 \(B\) 中的一个同态，它把 \(\gamma_{1}(x)\) 映为 \((0, x)\) ，对任意 \(x \in E\)）。

\begin{enumerate}
\def\labelenumi{\alph{enumi})}
\setcounter{enumi}{4}
\item
  设 E 和 F 是两个 k-模，h 是 E 到 F 的一个同态。证明存在唯一的分次代数同态 \(\Gamma(h):\Gamma(\mathrm{E})\to\Gamma(\mathrm{F})\) 使得 \(\gamma_{p}\circ h=\Gamma(h)\circ\gamma_{p}\) 对所有 \(p\in N\) 成立。验证若 g 和 h 是 k-模同态且 \(g\circ h\) 有定义，则 \(\Gamma(g\circ h)=\Gamma(g)\circ\Gamma(h)\) 。
\item
  设 \(h: \mathbf{E} \to \mathbf{F}\) 是 \(k\) -模的一个满同态。证明 \(\Gamma(h)\) 是满射，且其核是 \(\Gamma(\mathbf{E})\) 中由元素 \(\gamma_p(x)\) （其中 \(p > 0\) 且 \(x \in \operatorname{Ker}(h)\)）生成的理想。
\item
  设 E 和 F 是两个 k-模。证明存在唯一的同态 \(\varphi\) ，由代数 \(\Gamma(E \times F)\) 到代数 \(\Gamma(E) \otimes_{k} \Gamma(F)\) ，使得
\end{enumerate}

\[
(\varphi \circ \gamma_ {p}) (x, y) = \sum_ {r + s = p} \gamma_ {r} (x) \otimes \gamma_ {s} (y)
\]

对任意 \(p \in \mathbf{N}\) ， \(x \in E\) 和 \(y \in F\) 成立。证明这个同态是与分次相容的同构。

\begin{enumerate}
\def\labelenumi{\arabic{enumi})}
\setcounter{enumi}{2}
\tightlist
\item
  设 \(k\) 为一个交换环， \(E\) 为 \(k\) -模。证明存在唯一的同态 \(\Delta\) ，由代数 \(\Gamma(E)\) （练习 2）到代数 \(\Gamma(E) \otimes_k \Gamma(E)\) ，使得
\end{enumerate}

\[
(\Delta \circ \gamma_ {p}) (x) = \sum_ {r + s = p} \gamma_ {r} (x) \otimes \gamma_ {s} (x)
\]

对所有 \(x \in E\) 以及所有 \(p \in N\) 成立。证明余积 \(\Delta\) 是余结合的、余交换的（III，第 580-581 页），并在 \(\Gamma(E)\) 上定义了分次双代数结构（III，第 585 页）。证明在练习 2 c) 中定义的同态 \(\varepsilon\) 是一个余单位。

\begin{enumerate}
\def\labelenumi{\arabic{enumi})}
\setcounter{enumi}{3}
\item
  设 \(E_{1}\) 和 \(E_{2}\) 是交换环 k 上的两个模，且设 h 是 \(E_{1}\) 到 \(E_{2}\) 中的一个同态。验证 \(\Gamma(h)\) （练习 2，e））是分次双代数的同态（参见练习 3）。
\item
  设 \(\Gamma(\mathbf{Z})\) 是 \(\mathbf{Z}\) -模 \(\mathbf{Z}\) 的伽马代数（练习 2）。对每个 \(p \in \mathbf{N}\) ，令 \(\mathrm{T}^{(p)} = \gamma_p(1)\) 。
\end{enumerate}

\begin{enumerate}
\def\labelenumi{\alph{enumi})}
\tightlist
\item
  证明 \((\mathbf{T}^{(p)})_{p\in \mathbb{N}}\) 是 \(\mathbf{Z}\) -模 \(\Gamma (\mathbf{Z})\) 的一组基。证明存在唯一的同态 \(\theta\) ，由代数 \(\Gamma (\mathbf{Z})\) 到代数 \(\Gamma (\mathbf{Z})\otimes \Gamma (\mathbf{Z})\) ，使得
\end{enumerate}

\[
\theta \left(\mathrm{T} ^ {(p)}\right) = p! \mathrm{T} ^ {(p)} \otimes \mathrm{T} ^ {(p)}
\]

对所有 \(p \in \mathbf{N}\) 成立。

\begin{enumerate}
\def\labelenumi{\alph{enumi})}
\setcounter{enumi}{1}
\item
  设 B 是一个交换环，H 是环 \(\Gamma(\mathbf{Z})\) 到环 B 的同态组成的集合。证明对每个 \(f \in H\)，序列 \(f(\mathrm{T}^{(p)})\) （\(p \in N\)）是 B 中元素的一个指数型序列。证明映射 \(\varphi: f \mapsto (f(\mathrm{T}^{(p)}))_{p \in \mathbb{N}}\) 是 H 到 \(\mathcal{E}(\mathbf{B})\) 上的双射（练习 1）。
\item
  对 \(f, g \in \mathbf{H}\)，用 \(f * g\) 表示 \(\Gamma(\mathbf{Z}) \otimes \Gamma(\mathbf{Z})\) 到 \(\mathbf{B}\) 中的同态，它由 \((f * g)(u \otimes v) = f(u)g(v)\) （对任意 \(u, v \in \Gamma(\mathbf{Z})\)）定义。验证对任意 \(f, g \in \mathbf{H}\) 我们有
\end{enumerate}

\[
\begin{array}{c} \varphi (f) + \varphi (g) = \varphi ((f * g) \circ \Delta), \\ \varphi (f) \varphi (g) = \varphi ((f * g) \circ \theta). \end{array}
\]

\begin{enumerate}
\def\labelenumi{\arabic{enumi})}
\setcounter{enumi}{5}
\item
  设 \(n\) 为整数 \(> 0\)，并令 \(\Gamma(\mathbf{Z}/n\mathbf{Z})\) 为 \(\mathbf{Z}\) -模 \(\mathbf{Z}/n\mathbf{Z}\) 的伽马代数（练习 2）。证明对每个 \(p > 0\) ， \(\Gamma_p(\mathbf{Z}/n\mathbf{Z})\) 同构于 \(\mathbf{Z}\) 关于由整数 \(n^k((k, p - k))\) （其中 \(1 \leqslant k \leqslant p\)）所生成的子群所成的商群。证明若 \(p\) 与 \(n\) 是不同的素数，则 \(\Gamma_p(\mathbf{Z}/n\mathbf{Z})\) 是 \(n\) 阶循环群。证明若 \(n\) 是偶数，则 \(\Gamma_2(\mathbf{Z}/n\mathbf{Z})\) 是 \(2n\) 阶循环群。
\item
  （标量扩张）设 \(k\) 是一个交换环， \(L\) 为一个结合的交换含幺 \(k\) -代数， \(E\) 为 \(k\) -模。
\end{enumerate}

\begin{enumerate}
\def\labelenumi{\alph{enumi})}
\tightlist
\item
  设 \(\Gamma(L \otimes_k E)\) 为 L-模 \(L \otimes_k E\) 的伽马代数（练习 2）。证明存在唯一的同态 \(\theta\) ，由 \(L \otimes_k \Gamma(E)\) 到 \(\Gamma(L \otimes_k E)\) 中，与 \(L\) 上的代数结构相容，并且满足
\end{enumerate}

\[
\theta (1 \otimes \gamma_ {p} (x)) = \gamma_ {p} (1 \otimes x)
\]

对所有 \(x \in \mathbf{E}\) 和 \(p \in \mathbf{N}\) 成立。

\begin{enumerate}
\def\labelenumi{\alph{enumi})}
\setcounter{enumi}{1}
\tightlist
\item
  证明 \(\theta\) 是 \(L\) 上的分次双代数同构。
\end{enumerate}

\begin{enumerate}
\def\labelenumi{\arabic{enumi})}
\setcounter{enumi}{7}
\tightlist
\item
  设 \(k\) 为一个交换环， \(E\) 为 \(k\) -模。
\end{enumerate}

\begin{enumerate}
\def\labelenumi{\alph{enumi})}
\tightlist
\item
  证明存在唯一的代数 \(\Gamma(\mathrm{E})\) （练习 2）到 E 上的对称张量代数 TS(E)（IV，第 43 页）中的同态 g，使得
\end{enumerate}

\[
(g \circ \gamma_ {p}) (x) = x \otimes x \otimes \dots \otimes x \quad (p \text {   factors   } x),
\]

对任意 \(x \in \mathbb{E}\) ， \(p \in \mathbb{N}\) 成立。证明若 \(\mathbb{E}\) 是自由 \(k\) -模，以 \((e_i)_{i \in I}\) 为基，则 \(g\) 是分次双代数同构（参见练习 3 及 IV，第 51 页），且元素 \(\gamma_\alpha = \prod_i \gamma_{\alpha(i)}(e_i)\) （其中 \(\alpha \in \mathbf{N}^{(I)}\)）构成 \(\Gamma(\mathbb{E})\) 的一组基。

\begin{enumerate}
\def\labelenumi{\alph{enumi})}
\setcounter{enumi}{1}
\tightlist
\item
  设 \(\xi\) 是 \((1, n)\) 到 \(\mathbf{N}\) 中的一个映射，并令 \(p = \sum_{i=1}^{n} \xi(i)\) 。证明对 E 中元素的每个序列
\end{enumerate}

\(x_{1}, x_{2}, \ldots, x_{n}\)，\(\gamma_{\xi(1)}(x_{1}) \ldots \gamma_{\xi(n)}(x_{n})\) 在 \(g\) 下的像是 \(\sum_{\varphi} x_{\varphi(1)} \otimes \ldots \otimes x_{\varphi(p)}\)，其中 \(\varphi\) 遍历 \([1, p]\) 到 \([1, n]\) 中的所有映射组成的集合，使得

Card \(\varphi^{-1}(i) = \xi(i)\) 对每个 \(i \in [1, n]\) 成立（参见 IV，第 45 页，命题 3）。

\begin{enumerate}
\def\labelenumi{\alph{enumi})}
\setcounter{enumi}{2}
\tightlist
\item
  设 \(f\) 是 \(\otimes^p \mathbf{E}\) 到 \(\Gamma_p(\mathbf{E})\) 中的映射，由 \(f(x_1 \otimes \ldots \otimes x_p) = \gamma_1(x_1) \ldots \gamma_1(x_p)\) （对任意 \(x_1, \ldots, x_p \in \mathbf{E}\)）定义。证明对每个张量 \(t \in \otimes^p \mathbf{E}\) ， \((g \circ f)(t)\) 是 \(t\) 的对称化。证明对每个 \(u \in \Gamma_p(\mathbf{E})\) 我们有 \((f \circ g)(u) = p! u\) 。
\end{enumerate}

\begin{enumerate}
\def\labelenumi{\arabic{enumi})}
\setcounter{enumi}{8}
\tightlist
\item
  （多项式律）用 E 和 F 表示交换环 k 上的两个模。所谓从 E 到 F 的多项式律，我们理解为：对每个（结合、交换且含单位元的）k-代数 L，一个映射 \(\varphi_{L}\)，由 \(L \otimes_{k} E\) 到 \(L \otimes_{k} F\) 中，满足条件：
\end{enumerate}

若 L 和 R 是两个 k-代数且 f 是 L 到 R 的保单位同态，则

\[
(f \otimes 1) \circ \varphi_ {L} = \varphi_ {R} \circ (f \otimes 1).
\]

\begin{enumerate}
\def\labelenumi{\alph{enumi})}
\tightlist
\item
  设 \(\varphi\) 是从 \(E\) 到 \(F\) 中的一个多项式律，设 \(L = k[(T_i)_{i \in I}]\) 为一族未定元 \(T_i\) 在 \(k\) 上的多项式代数，并令 \((x_i)_{i \in I}\) 是 \(E\) 的元素的一个有限支撑族。证明存在唯一的有限支撑族 \((y_\xi)_{\xi \in N^{(I)}}\) （由 \(F\) 的元素组成），使得
\end{enumerate}

\[
\varphi_ {L} \left(\sum_ {i} T _ {i} \otimes x _ {i}\right) = \sum_ {\xi \in N ^ {(I)}} T ^ {\xi} \otimes y _ {\xi}.
\]

证明若 \(\mathbf{R}\) 是 \(k\) 上的代数，则对每个族 \((a_i)_{i \in I}\) （元素取自 \(\mathbf{R}\) ），我们有

\[
\varphi_ {R} \left(\sum_ {i} a _ {i} \otimes x _ {i}\right) = \sum_ {\xi \in \mathbf {N} ^ {(I)}} a ^ {\xi} \otimes y _ {\xi}.\tag{1}
\]

\begin{enumerate}
\def\labelenumi{\alph{enumi})}
\setcounter{enumi}{1}
\tightlist
\item
  E 到 F 的多项式律的族 \((\varphi^j)_{j \in \mathbb{J}}\) 称为可和的，如果对 \(k\) 上的每个代数 L 以及每个 \(u \in \mathrm{L} \otimes \mathrm{E}\)，族 \(\varphi_L^j(u)\) 具有有限支撑。证明若 \((\varphi^j)_{j \in \mathbb{J}}\) 是 E 到 F 的多项式律的可和族，则存在唯一的从 E 到 F 的多项式律 \(\varphi\)，使得对每个 \(k\) -代数 L 以及每个 \(u \in \mathrm{L} \otimes \mathrm{E}\) 我们有 \(\sum_j \varphi_L^j(u) = \varphi_L(u)\) 。这个多项式律称为诸律 \(\varphi^j\) 的和，记为
\end{enumerate}

\[
\sum_ {j} \varphi^ {j}.
\]

\begin{enumerate}
\def\labelenumi{\alph{enumi})}
\setcounter{enumi}{2}
\item
  设 E、F、G 为三个 k-模， \(\varphi\) 为从 E 到 F 的一个多项式律， \(\psi\) 为从 F 到 G 的一个多项式律。证明存在唯一的从 E 到 G 的多项式律 \(\eta\)，使得对每个 k-代数 L 都有 \(\eta_{L} = \psi_{L} \circ \varphi_{L}\) 。律 \(\eta\) 称为 \(\varphi\) 与 \(\psi\) 的复合，记为 \(\psi \circ \varphi\) 。
\item
  设 \(p\) 为整数 \(\geqslant 0\)。多项式律 \(\varphi\) （从 \(E\) 到 \(F\) 中）称为 \(p\) 次齐次的，如果对每个代数 \(L\) （在 \(k\) 上）我们有 \(\varphi_{L}(au) = a^{p}\varphi_{L}(u)\) （对任意 \(a \in L\) 和 \(u \in L \otimes E\)）。证明若 \(\varphi\) 是 \(p\) 次齐次的，则在公式 (1) 中我们有 \(y_{\xi} = 0\) （对 \(|\xi| = \sum_{i} \xi(i) \neq p\)）。确定所有 0 次和 1 次的齐次多项式律。
\item
  使用与 c) 中相同的记号，证明若 \(\varphi\) 与 \(\psi\) 分别是 \(p\) 次和 \(q\) 次的齐次律，则 \(\psi \circ \varphi\) 是 \(pq\) 次的齐次律。
\item
  设 \(p\) 为整数 \(\geqslant 0\) 。对每个交换 \(k\) -代数 \(L\)，我们用 \(\varepsilon_L\) 表示单射 \(a \mapsto aT\) （从 \(L\) 到多项式代数 \(L[T]\) 中），并用 \(\beta_L^p\) 表示 \(L\) -模 \(L[T]\) 到 \(L\) 中的同态，它把每个元素 \(\sum_k a_k T^k \in L[T]\) 对应到系数 \(a_p\) 。设 \(\varphi\) 是从 \(E\) 到 \(F\) 中的一个多项式律。证明
\end{enumerate}

映射 \(\varphi_{L}^{p} = (\beta_{L}^{p}\otimes \mathrm{Id}_{F})\circ \varphi_{L[T]}\circ (\varepsilon_{L}\otimes \mathrm{Id}_{E})\) 构成从 E 到 F 的一个多项式律，它是 \(p\) 次齐次的。这个多项式律称为 \(p\) 次齐次分量，对应于律 \(\varphi\) 。证明 \(\varphi\) 的齐次分量构成一个可和族，其和为 \(\varphi\) 。

\begin{enumerate}
\def\labelenumi{\arabic{enumi})}
\setcounter{enumi}{9}
\tightlist
\item
  记 \(k\) 为一个交换环，\(E\) 为一个 \(k\) -模。对每个整数 \(p \geqslant 0\) 以及每个 \(k\) -代数 \(L\)，我们用 \(\gamma_{p,L}\) 表示 \(L \otimes E\) 到 \(L \otimes \Gamma_p(E)\) 中的映射，它由 \(\gamma_p: L \otimes E \to \Gamma_p(L \otimes E)\) 与 \(\Gamma_p(L \otimes E)\) 到 \(L \otimes \Gamma_p(E)\) 上的典范同构复合而成（参见练习 2 和练习 7）。
\end{enumerate}

\begin{enumerate}
\def\labelenumi{\alph{enumi})}
\tightlist
\item
  验证映射 \(\gamma_{p,L}\) 构成 E 到 \(\Gamma_p(E)\) 中的一个多项式律（练习 9），它是 \(p\) 次齐次的；这条律在下文中将记作 \(\boldsymbol{\gamma}_p\) 。证明对每个 \(k\) -模 F 以及每个 \(\sigma \in \operatorname{Hom}(\Gamma_p(E), F)\)，映射 \(\varphi_L = (\operatorname{Id}_L \otimes \sigma) \circ \boldsymbol{\gamma}_{p,L}\) 构成 E 到 F 的一个多项式律，它是 \(p\) 次齐次的。这条律称为与 \(\sigma\) 相伴的律。
\end{enumerate}

设 \((x_{i})_{i\in I}\) 是 E 的一族元素，且 \((a_{i})_{i\in I}\) 是 L 中具有有限支撑的一族元素。证明

\[
\varphi_ {L} \left(\sum_ {i} a _ {i} \otimes x _ {i}\right) = \sum_ {\xi \in N ^ {(I)}, | \xi | = p} a ^ {\xi} \otimes \sigma (\gamma_ {\xi} (x))
\]

，其中 \(|\xi| = \sum_{i} \xi(i), a^{\xi} = \prod_{i} a_{i}^{\xi(i)}\) 且 \(\gamma_{\xi}(x) = \prod_{i} \gamma_{\xi(i)}(x_{i})\) 。

\begin{enumerate}
\def\labelenumi{\alph{enumi})}
\setcounter{enumi}{1}
\tightlist
\item
  设 \(\varphi\) 是 E 到 k-模 F 中的 p 次齐次多项式律。证明存在唯一的线性映射 \(\sigma\)（从 \(\Gamma_{p}(\mathrm{E})\) 到 F 中），使得 \(\varphi\) 是与 \(\sigma\) 相伴的律。（首先假设 E 以 \((x_i)_{i \in I}\) 为基；利用练习 9，证明在此情形存在唯一的映射 \(\sigma \in \operatorname{Hom}(\Gamma_p(E), F)\)，使得对每个代数 L 及每个由 L 中元素组成的有限支撑族 \((a_i)_{i \in I}\)，练习 9 的关系式 (1) 成立。通过引入一个自由 \(k\) -模 \(E'\) 和一个满同态 \(h: E' \to E\) 来处理一般情形。多项式律 \(\varphi\) 与由 \(h\) 定义的律的复合是 \(E'\) 到 F 的一个多项式律，它相伴于一个线性映射 \(\sigma': \Gamma_p(E') \to F\)；证明 \(\sigma'\) 是 \(\Gamma_p(h)\) 与一个映射 \(\sigma: \Gamma_p(E) \to F\) 的复合，且 \(\varphi\) 相伴于 \(\sigma\)（参见练习 2， \(f\)）。）
\end{enumerate}

\begin{enumerate}
\def\labelenumi{\arabic{enumi})}
\setcounter{enumi}{10}
\tightlist
\item
  我们沿用练习 10 的记号。\(\Gamma(E)\) 到 \(F\) 中的线性映射将称为 \(E\) 上取值于 \(F\) 的级数。\(E\) 上取值于 \(F\) 的级数构成一个 \(k\) -模，记作 \(S(E, F)\)。对 \(n \in N\)，我们说级数 \(\sigma \in S(E, F)\) 的阶为 \(\geqslant n\)，如果 \(\operatorname{Ker}(\sigma) \supset \Gamma_p(E)\) 对所有 \(p < n\) 成立。在下文中，我们将赋予 \(S(E, F)\) 拓扑群结构，以诸 \(k\) -模（由阶 \(\geqslant n\) 的级数组成，其中 \(n \in N\)）为 0 的邻域基。
\end{enumerate}

设 \(\varphi\) 是从 \(E\) 到 \(F\) 中的一个多项式律。证明存在唯一的级数 \(\sigma \in S(E, F)\)，使得：

\begin{enumerate}
\def\labelenumi{(\roman{enumi})}
\tightlist
\item
  对每个 \(k\) -代数 \(\mathbf{L}\) 以及每个 \(u \in \mathbf{L} \otimes \mathbf{E}\)，族 \((\mathrm{Id}_{\mathbf{L}} \otimes \sigma) \gamma_{p,\mathbf{L}}(u) (p \in \mathbf{N})\) 具有有限支撑；
\end{enumerate}

\[
\varphi_ {L} (u) = \sum_ {p} \left(\operatorname{Id} _ {L} \otimes \sigma\right) \circ \gamma_ {p, L} (u). \tag {ii}
\]

（应用练习 10，b) 的结果）。

级数 \(\sigma\) 称为与律 \(\varphi\) 相伴的级数。

\begin{enumerate}
\def\labelenumi{\arabic{enumi})}
\setcounter{enumi}{11}
\tightlist
\item
  设 \(k\) 是一个交换环， \(E\) 为 \(k\) -模， \(B\) 为 \(k\) -代数（不一定是结合的）。设 \(\Delta\) 为 \(\Gamma(E)\) 的余积（参见练习 3）。对任意级数 \(\sigma, \sigma' \in S(E, B)\)，我们用 \(\sigma \cdot \sigma'\) 表示级数 \(m \circ (\sigma \otimes \sigma') \circ \Delta\)，其中 \(m\) 表示 \(B \otimes_k B\) 到 \(B\) 中的、由 \(B\) 中乘法定义的线性映射。
\end{enumerate}

\begin{enumerate}
\def\labelenumi{\alph{enumi})}
\item
  证明映射 \((\sigma, \sigma') \mapsto \sigma \cdot \sigma'\) 在 S(E, B) 上定义了一个 \(k\) -代数结构（参见练习 11）。证明若 B 是结合的（相应地，交换的、含单位元的），则 S(E, B) 是结合的（相应地，交换的、含单位元的）（参见 III，第 578-584 页）。
\item
  假设 B 是结合的且含单位元。证明 S(E, B) 中的可逆元素是那些使 σ(1) 在 B 中可逆的级数 σ。
\item
  假设 E 是一个自由 \(k\) -模，并且 \((e_i)_{i \in I}\) 是 E 的一组基。对每个 \(\alpha \in \mathbf{N}^{(I)}\)，令 \(e^{(\alpha)} = \Pi \gamma_{\alpha(i)}(e_i)\)，并用 \(f^\alpha\) 表示 E 上取值于 \(k\) 的、由 \(f^\alpha(e^{(\beta)}) = \delta_{\alpha, \beta}\)（对每个 \(\beta \in \mathbf{N}^{(I)}\)，克罗内克指标）定义的级数（参见练习 8）。证明 \(f^\alpha f^\beta = f^{\alpha + \beta}\) 对任意 \(\alpha, \beta \in \mathbf{N}^{(I)}\) 成立。由此推出，当 I 有限时，代数 \(S(E, k)\) 同构于形式幂级数代数 \(k[[(X_i)_{i \in I}]\) 。
\end{enumerate}

\begin{enumerate}
\def\labelenumi{\arabic{enumi})}
\setcounter{enumi}{12}
\tightlist
\item
  （级数的分幂。）记号同练习 10 和练习 11。设 E 和 F 为两个 \(k\) -模，\(\sigma\) 为 E 上取值于 F 的一个级数。用 \(\sigma\) 表示 \(\Gamma(E)\) 到 F 中、由 \(\sigma\) 所定义的 1 次齐次多项式律。对任意整数 \(p\) , \(m \in \mathbf{N}\)，我们用 \(\sigma_m^{(p)}\) 表示与多项式律 \(\gamma_p \circ \sigma \circ (\gamma_0 + \gamma_1 + \ldots + \gamma_m)\) 相伴的级数（参见练习 10）。
\end{enumerate}

\begin{enumerate}
\def\labelenumi{\alph{enumi})}
\item
  证明当 \(m\) 趋于 \(\infty\) 时，序列 \(\sigma_{m}^{(p)}\) 在 \(S(E, \Gamma_p(F))\) 中收敛于一个级数 \(\sigma^{(p)}\) 。级数 \(\sigma^{(p)}\) 称为分 \(p\) 次幂（对级数 \(\sigma\) 而言）。
\item
  证明如果 \((x_{i})_{i\in I}\) 是 \(E\) 中元素的一个族，且 \(\eta \in \mathbf{N}^{(I)}\) ，则
\end{enumerate}

\[
\sigma^ {(p)} \left(\gamma_ {\eta} (x)\right) = \sum_ {\xi \in \mathscr {L}} \prod_ {\alpha \in N ^ {(I)}} \gamma_ {\xi (\alpha)} \left(\sigma \left(\gamma_ {\alpha} (x)\right)\right),
\]

，其中 \(\mathcal{L}\) 是 \(\mathbf{N}^{(I)}\) 到 \(\mathbf{N}\) 中的、满足条件 \(\sum_{\alpha} \xi(\alpha) = p\) 与 \(\sum_{\alpha} \xi(\alpha) \alpha = \eta\) 的所有有限支撑映射的集合（应用练习 10， \(a\)）。

\begin{enumerate}
\def\labelenumi{\alph{enumi})}
\setcounter{enumi}{2}
\tightlist
\item
  证明如果 \(\sigma, \tau \in S(E, F)\)，则对每个 \(p \in \mathbf{N}\) 我们有
\end{enumerate}

\[
\left(\sigma + \tau\right) ^ {(p)} = \sum_ {r + s = p} \sigma^ {(r)} \sigma^ {(s)}
\]

（乘积在代数 \(S(E, \Gamma(E))\) 中计算，参见练习 12）。

\(d\) ）证明对每个整数 \(p\in \mathbf{N}\) 我们有 \(\sigma^p = p!\sigma^{(p)}\) ，其中 \(\sigma^p\) 是 \(p\) 次幂，即 \(\sigma\) 在代数 \(\mathrm{S}(E, \Gamma(F))\) 中的幂

\begin{enumerate}
\def\labelenumi{\alph{enumi})}
\setcounter{enumi}{4}
\tightlist
\item
  设 \(\Delta_{\mathrm{E}}\) 与 \(\Delta_{\mathrm{F}}\) 分别为 \(\Gamma(\mathrm{E})\) 与 \(\Gamma(\mathrm{F})\) 中的余积。证明对所有 \(p \in \mathbf{N}\) 我们有
\end{enumerate}

\[
\Delta_ {F} \circ \sigma^ {(p)} = \sum_ {r + s = p} \left(\sigma^ {(r)} \otimes \sigma^ {(s)}\right) \circ \Delta_ {E}.
\]

\begin{enumerate}
\def\labelenumi{\arabic{enumi})}
\setcounter{enumi}{13}
\tightlist
\item
  沿用练习 13 中的记号，假设 \(\sigma \in S(E, F)\) 的阶为 \(\geqslant 1\) 。
\end{enumerate}

\begin{enumerate}
\def\labelenumi{\alph{enumi})}
\tightlist
\item
  证明对每个 \(p \in \mathbf{N}\)，\(\sigma^{(p)}\) 的阶为 \(\geqslant p\) 。由此推出族 \((\sigma^{(p)})_{p \in \mathbf{N}}\) 是可和的。记 \(\mathbf{A}(\sigma) = \sum_{p} \sigma^{(p)}\)，验证对每个 \(n \in \mathbf{N}\)，
\end{enumerate}

\[
\mathbf {A} (\sigma) \left(\Gamma_ {n} (\mathrm{E})\right) \subset \Gamma_ {0} (\mathrm{F}) + \dots + \Gamma_ {n} (\mathrm{F}).
\]

\begin{enumerate}
\def\labelenumi{\alph{enumi})}
\setcounter{enumi}{1}
\item
  证明如果 \(\operatorname{Ker}(\sigma) \supset \Gamma_m(E)\) 对所有 \(m \neq 1\) 成立，则 \(A(\sigma) = \Gamma(\sigma \circ \gamma_1)\) 。
\item
  证明 \(\mathbf{A}(\sigma)\) 是余代数 \(\Gamma(\mathbf{E})\) 到余代数 \(\Gamma(\mathbf{F})\) 中的同态，且与余单位相容（参见练习 13，e)）。
\item
  证明如果 \(\sigma, \tau \in S(E, F)\) 是两个阶为 \(\geqslant 1\) 的级数，则 \(A(\sigma + \tau) = A(\sigma)A(\tau)\)，其中乘积 \(A(\sigma)A(\tau)\) 由 \(\Gamma(E)\) 中的乘积定义（参见练习 13，c)）。
\end{enumerate}

\begin{enumerate}
\def\labelenumi{\arabic{enumi})}
\setcounter{enumi}{14}
\tightlist
\item
  沿用练习 14 中的记号，设 E、F、G 为三个 k-模， \(\sigma \in S(E, F)\)，\(\tau \in S(F, G)\)，并假设 \(\sigma\) 的阶为 \(\geqslant 1\) 。级数 \(\tau \circ A(\sigma) \in S(E, G)\) 称为 \(\sigma\) 与 \(\tau\) 的复合；记作 \(\tau \circ \sigma\) 。
\end{enumerate}

\begin{enumerate}
\def\labelenumi{\alph{enumi})}
\item
  设 \(\varphi\) 是 E 到 F 的一个多项式律（练习 9），\(\psi\) 是 F 到 G 的一个多项式律，并记 \(\sigma_{\varphi}\) 与 \(\sigma_{\psi}\) 分别为与 \(\varphi\) 和 \(\psi\) 相伴的级数。证明如果 \(\varphi\) 的 0 次齐次分量为零，则 \(\sigma_{\varphi}\) 的阶为 \(\geqslant 1\)，且 \(\sigma_{\psi} \circ \sigma_{\varphi}\) 是与多项式律 \(\psi \circ \varphi\) 相伴的级数。
\item
  证明如果 \(\sigma\) 的阶为 \(\geqslant p\) 且 \(\tau\) 的阶为 \(\geqslant q\) ，则 \(\tau \circ \sigma\) 的阶为 \(\geqslant pq\) 。
\item
  证明对每个整数 \(p\)，\((\tau \circ \sigma)^{(p)} = \tau^{(p)} \circ A(\sigma)\) 。由此推出，如果 \(\tau\) 的阶为 \(\geqslant 1\) 且 \(H\) 是一个 \(k\) -模，则对任意 \(\rho \in S(G, H)\) 我们有 \((\rho \circ \tau) \circ \sigma = \rho \circ (\tau \circ \sigma)\) 。 d) 假设 \(G\) 是一个 \(k\) -代数。证明映射 \(\tau \mapsto \tau \circ A(\sigma)\) 是代数 \(S(F, G)\) 到代数 \(S(E, G)\) 中的同态。
\end{enumerate}

\begin{enumerate}
\def\labelenumi{\arabic{enumi})}
\setcounter{enumi}{15}
\tightlist
\item
  设 E 和 F 是交换环 \(k\) 上的两个模，并令 \(p\) 为一个整数 \(\geqslant 0\) 。
\end{enumerate}

\begin{enumerate}
\def\labelenumi{\alph{enumi})}
\item
  证明如果 \(p \leqslant 2\)，映射 \(\lambda_p\)——即 \(\operatorname{Hom}(\Gamma_p(E), F)\) 到 \(F^E\) 中的、把 \(\sigma\) 对应到映射 \(\sigma \circ \gamma_p\) 的映射——是单射。
\item
  证明 \(\lambda_1\) 以 \(\operatorname{Hom}_k(E, F)\) 为像。证明 \(\lambda_2\) 的像是所有满足下列条件的映射 \(g\) （从 \(E\) 到 \(F\) 中）的集合：
\end{enumerate}

\begin{enumerate}
\def\labelenumi{(\roman{enumi})}
\item
  \(g(ax) = a^2 g(x)\) 对任意 \(a \in k\) 与 \(x \in E\) 成立；
\item
  映射 \((x, y) \mapsto g(x + y) - g(x) - g(y)\) 是 \(E \times E\) 到 \(F\) 中的一个双线性映射。c) 证明如果 \(E\) 是一个自由 \(k\) -模，则 \(\lambda_p\) 的像是模 \(\operatorname{Pol}_k^p(E, F)\) ，即 \(p\) 次齐次多项式映射（从 \(E\) 到 \(F\) 中）所成的模。
\end{enumerate}

